\subsection{A positive summand form and an explicit tail bound}

The combination in \eqref{gaussian:eq:cubic-compact} can be computed without cancellation
between separate Tornheim constants.  For $k=m+n$, define
\[
 R_A(m,n)=A^2+2A\log k+\log m\log k+\log n\log k-\log m\log n.
\]
Since $A>0$, $\log k\ge\log m,\log n\ge0$, we have $R_A(m,n)>0$.
Writing $p=\log m,q=\log n,r=\log k$, one also has
$r(p+q)-pq\le r^2$: for fixed $r$, the expression is increasing separately
in $p,q\in[0,r]$.  Consequently
\[
 0<R_A(m,n)\le(A+\log k)^2.
\]
Therefore the diagonal partial sums
\[
 C_N=\frac3{32}\sum_{m+n\le N}\frac1{mn(m+n)}
    +\frac{3}{8\pi^2}
       \sum_{m+n\le N}\frac{R_A(m,n)}{mn(m+n)}
\]
increase to the centered cubic moment $M_3-a^3-3aV$.  Moreover,
\begin{align*}
 0<(M_3-a^3-3aV)-C_N
 &\le\sum_{k>N}\frac{2(1+\log k)}{k^2}
       \left(\frac3{32}+\frac{3(A+\log k)^2}{8\pi^2}\right).
\end{align*}
This gives a transparent computable enclosure, although its convergence
is only $O(\log^3N/N)$.  An integral bound may be used once the explicit
summand is decreasing; direct summation can handle the finite initial range.
This positivity observation is an immediate corollary of the Fourier
calculation, not a claim of a new classical-constant evaluation.

\subsection{Exponentially convergent numerical reproduction}

For $t>0$, put
\[
 C(t)=\operatorname{Li}_1(e^{-t})=-\log(1-e^{-t}),\qquad
 F(t)=\left.\partial_s\operatorname{Li}_s(e^{-t})\right|_{s=1}.
\]
The usual Mellin representation of the Tornheim sum gives
\begin{equation}\label{tornheim:cubic:mellinpartials}
 T_{12}=\int_0^\infty F(t)^2\,dt,\quad
 T_{13}=\int_0^\infty F(t)C(t)(\log t+\gamma)\,dt,\quad
 T_3=\int_0^\infty C(t)^2(\log t+\gamma)\,dt.
\end{equation}
For $0<t<2\pi$, cancellation of the two poles in Jonqui\`ere's expansion
produces
\begin{align*}
 C(t)&=-\log t+\sum_{k\ge1}c_kt^k,
 &c_k&=\frac{(-1)^k\zeta(1-k)}{k!},\\
 F(t)&=-\alpha-\gamma\log t-\tfrac12\log^2t+\sum_{k\ge1}b_kt^k,
 &b_k&=\frac{(-1)^k\zeta'(1-k)}{k!},\\
 \alpha&=\gamma_1+\frac{\gamma^2}2+\frac{\pi^2}{12}.
\end{align*}
Here the Stieltjes convention is
$\zeta(1+z)=z^{-1}+\sum_{j\ge0}(-1)^j\gamma_jz^j/j!$.
For computational stability the coefficient formulas can be rewritten as
\begin{align*}
 c_1&=\tfrac12,&b_1&=\tfrac12\log(2\pi),\\
 c_{2j}&=\frac{(-1)^j\zeta(2j)}{j(2\pi)^{2j}},&
 b_{2j}&=c_{2j}\left(\log(2\pi)-\psi(2j)-\frac{\zeta'(2j)}{\zeta(2j)}\right),\\
 c_{2j+1}&=0\quad(j\ge1),&
 b_{2j+1}&=\frac{(-1)^{j+1}\pi\zeta(2j+1)}{(2j+1)(2\pi)^{2j+1}}
 \quad(j\ge1).
\end{align*}
The functional equation of the Riemann zeta function proves these identities.
Elementary bounds $\zeta(k)<2$, $|\zeta'(k)|<2$, $H_{k-1}\le k-1$,
$\gamma<1$, $\log(2\pi)<2$, and $2\pi>6$ give, for $k\ge2$,
\[
 |b_k|\le12\,6^{-k},\qquad |c_k|\le2\,6^{-k}.
\]
Thus after degree $K\ge1$ their remainders on $0<t\le1$ are at most
\[
 \epsilon_b=\frac{12}{5}6^{-K},\qquad
 \epsilon_c=\frac25 6^{-K}.
\]
All integrals of the resulting polynomials in $t$ and $\log t$ are evaluated
using $\int_0^1t^k\log^jt\,dt=(-1)^jj!/(k+1)^{j+1}$.

On $t\ge1$, insert
$C(t)=\sum_{n\ge1}e^{-nt}/n$ and
$F(t)=-\sum_{n\ge1}(\log n)e^{-nt}/n$.  The integral with the logarithmic
factor uses
\[
 \int_1^\infty e^{-kt}(\log t+\gamma)\,dt
 =\frac{E_1(k)+\gamma e^{-k}}k.
\]
Truncate this sum at $m+n\le N$.  A common upper bound for the three
omitted tails is
\[
 B_N=4\sum_{k>N}ke^{-k}
 =\frac{4e^{-(N+1)}((N+1)-Ne^{-1})}{(1-e^{-1})^2}.
\]
For the three integrals in the order $T_{12},T_{13},T_3$, safe total
series-truncation bounds are respectively
\[
 \epsilon_b(12+\epsilon_b)+B_N,\qquad
 5\epsilon_b+15\epsilon_c+2\epsilon_b\epsilon_c+B_N,\qquad
 10\epsilon_c+2\epsilon_c^2+B_N.
\]
For example the elementary estimates
$C(t)\le-\log t+1$ and
$|F(t)|\le4-\log t+\tfrac12\log^2t$ on $(0,1]$ imply all the constants
above.  The sharper bound $\int_0^1|F(t)|dt<6$ follows directly from its
coefficient expansion (one may use $0<\alpha<2$).
An elementary verification of the latter constant bound is available:
the defining limit for $\gamma_1$ and the total variation of
$x\mapsto\log x/x$ on $[1,\infty)$ give $|\gamma_1|\le2/e<3/4$;
combine this with $0<\gamma<3/5$ and $3<\pi<\sqrt{12}$.
This is an analytic truncation certificate; floating-point rounding must
be handled separately to turn an implementation into interval arithmetic.

The independent implementation \texttt{verify\_tornheim.py} evaluates the
three partials and compares \eqref{gaussian:eq:cubic-compact} against direct log-gamma
quadrature.  Its output explicitly labels the rounding qualification.

\subsection{A sign check for an auxiliary identity}

If the June 2012 preprint is used, an auxiliary sign should be checked
before copying its Remark~3.  Put
\[
 S_\gamma=\sum_{m,n\ge1}
        \frac{(\gamma+\log m)(\gamma+\log n)}{mn(m+n)}.
\]
Directly expanding the numerator, with absolutely convergent sums, gives
\[
 T_{12}=S_\gamma+2\gamma\bigl(\zeta'(2,1)+\zeta'(3)\bigr)
                  -2\gamma^2\zeta(3).
\]
The coefficient of $S_\gamma$ is positive.  The displayed minus sign in
equations~(63) and~(65) of the author-hosted June 2012 preprint would make
their right side negative, although $T_{12}>0$.  This issue does not
affect the cubic formula proved above.  This note does not assert that
the same typographical error occurs in the final journal typesetting.

